\section{Mechanical minimum patterns and spectral descent}
\label{sec:spectral}

When the product index is zero, the estimate used in the preceding
section loses its factor $\lambda^{-1}$ per lower-state visit. A
different input is needed. Cyclic balance controls each contiguous
block of lower states, and the equality itself turns this control into
an integer gap in the block exponent. We then measure the Vieta error
in the spectral coordinate $\arcosh(3L/2)$. In that coordinate a
mutation is almost a subtraction, and the integer coefficients of
the Euclidean algorithm control the accumulated error.

\begin{theorem}[The minimum-index parameter bound]
\label{thm:minimum-bound}
In the minimum-ray model \eqref{eq:gap-ray-data}--\eqref{eq:gap-positive-model},
let $h\geqslant2$ and let $\epsilon_1,\ldots,\epsilon_h$ be a cyclic
mechanical binary pattern with $1\leqslant t=\sum_i\epsilon_i<h$.
If
\begin{equation}\label{eq:spectral-minimum-equality}
 V_0=\frac13\tr\prod_{i=1}^h A_0B^{\epsilon_i}=u_{\eps m},
 \qquad m\geqslant1,
\end{equation}
then
\begin{equation}\label{eq:spectral-final-bound}
 p<(100h^2)^{h+1}.
\end{equation}
The analytic statement does not require $\gcd(h,t)=1$.
\end{theorem}

Here cyclic mechanical means that, for some intercept $\omega$, the
bits are successive differences of
$\lfloor it/h+\omega\rfloor$, extended periodically with period $h$.
For every proper cyclic interval of $j$ positions, its number $e$ of
ones therefore satisfies
\begin{equation}\label{eq:spectral-balance}
 |e-tj/h|<1.
\end{equation}
Indeed the difference from $tj/h$ is the difference of the fractional
parts at the two endpoints. This strict interval estimate is the only
mechanical-word property used in the proof.

\subsection{Choose a dominant state before estimating the error}

Set $\kappa=9K=1+e_0$, where $e_0=4/D$, as in
\eqref{eq:positive-kappa}. Choose the larger of the two constant-state
contributions to $V_0$. If it is the upper contribution, put $\gamma=1$
and $(c_0,d_0)=(c,d)$. Otherwise interchange the two matrix coordinates,
put $\gamma=-1$ and $(c_0,d_0)=(d,c)$, and use signed gaps
$\gamma\epsilon_i$ with the displayed diagonal matrix $B$.
This operation leaves the word trace unchanged. The represented ray
is now $u_i$ or $u_{-i}$ respectively, and the signed single index is
$\gamma\eps m$. In particular its magnitude and minimum normalization
are unchanged. These transformations are changes of the trace model;
no surface isometry is asserted.

Put
\begin{equation}\label{eq:spectral-dominant-data}
 \begin{gathered}
 \tau=\gamma t,\qquad
 X=3^{h-1}c_0^h\lambda^\tau,\qquad
 C=K\kappa^{h-1},\\
 \rho=d_0/c_0,\qquad Z=C/X^2\leqslant1,
 \qquad \sigma=4/(9p^2).
 \end{gathered}
\end{equation}
The other constant contribution is $C/X$. If $\mathcal M$ denotes
the sum of all mixed path contributions, then
\begin{equation}\label{eq:spectral-trace-split}
 V_0=X(1+Z)+\mathcal M.
\end{equation}
For a mixed path, the proper set of lower positions consists of $r$
nonempty contiguous cyclic blocks. Let $j$ be its total size and $e$
its total number of original one bits. The path expansion of
Section~\ref{sec:gap} gives the exact relative weight
\begin{equation}\label{eq:spectral-block-weight}
 w=\sigma^r\rho^j\lambda^{-2\gamma e}.
\end{equation}

For a block of size $j_i$ and one-count $e_i$, the number
$1+\gamma(e_i-tj_i/h)$ is positive by
\eqref{eq:spectral-balance}. It is a multiple of $1/h$, and hence
at least $1/h$. Since
$Z=\rho^h\lambda^{-2\gamma t}$, multiplying these block estimates
in \eqref{eq:spectral-block-weight} gives
\begin{equation}\label{eq:spectral-coarse-block-weight}
 w\leqslant a_0^r z_0^j,
 \qquad a_0=\sigma\lambda^2\lambda^{-2/h}<4\lambda^{-2/h},
 \qquad z_0=Z^{1/h}\leqslant1.
\end{equation}
The exponent gain is paid once per lower block, rather than once per
position. This distinction will also control the number of relevant paths.

\subsection{Sum paths by their transitions}

For positive scalars $a,z$, consider
$$
 T(a,z)=\begin{pmatrix}1&\sqrt{az}\\\sqrt{az}&z\end{pmatrix},
 \qquad \Phi_h(a,z)=\tr T(a,z)^h-1-z^h.
$$
A mixed cyclic state path with $j$ lower positions and $2r$ transitions
contributes exactly $a^rz^j$ to this trace. All coefficients are positive.
Choosing its $2r$ transition edges and an initial state gives exactly
$2\binom h{2r}$ such paths when the lower-state size is not prescribed.
It follows that
\begin{equation}\label{eq:spectral-partition-function}
 \begin{aligned}
 \Psi_h(a):=\Phi_h(a,1)
 &=2\sum_{r=1}^{\lfloor h/2\rfloor}\binom h{2r}a^r\\
 &=(1+\sqrt a)^h+(1-\sqrt a)^h-2.
 \end{aligned}
\end{equation}
For $h^2a<1$, the inequality $2\binom h{2r}\leqslant h^{2r}$
gives the useful rational bound
\begin{equation}\label{eq:spectral-partition-bound}
 \Psi_h(a)\leqslant\frac{h^2a}{1-h^2a}.
\end{equation}
Thus \eqref{eq:spectral-coarse-block-weight} gives, before any
small-discrepancy assumption,
\begin{equation}\label{eq:spectral-first-path-bound}
 \mathcal M/X\leqslant\Phi_h(a_0,z_0)\leqslant\Psi_h(a_0).
\end{equation}
This is a sum of positive indexed paths, so cyclic repetitions or a
nonprimitive period do not change its validity.

Suppose henceforth, for a contradiction, that
\begin{equation}\label{eq:spectral-threshold-assumption}
 p\geqslant S^{h+1},\qquad S=100h^2.
\end{equation}
Then $\lambda>S^{h+1}$ and $h^2a_0<4h^2/S^2<1/16$.
Equations \eqref{eq:spectral-partition-bound} and
\eqref{eq:spectral-first-path-bound} imply
\begin{equation}\label{eq:spectral-coarse-mixed}
 \mathcal M/X<\frac{64h^2}{15S^2}<\frac1{16}.
\end{equation}
Let $a$ be the leading coefficient of the single entry in the
transformed ray, so $a$ is $c_0$ or $d_0$, and put $Y=a\lambda^m$.
Lemma~\ref{lem:positive-discrepancy} gives $Y>X$.
The minimum ratio also gives $Y>\sqrt\lambda/3$.
By \eqref{eq:spectral-trace-split}, $Z\leqslant1$ and
\eqref{eq:spectral-coarse-mixed}, we have $V_0<3X$; hence
$X>\sqrt\lambda/9$. Since $C<2^h$ and $S\geqslant400>32$,
\begin{equation}\label{eq:spectral-opposite-constant}
 Z<\frac{81\cdot2^h}{\lambda}
 <\frac{81}{32\cdot16^h}\leqslant\frac{81}{8192}<\frac1{16}.
\end{equation}
The exact equality $Y+K/Y=V_0$ now shows that
\begin{equation}\label{eq:spectral-small-delta}
 0<\delta:=Y/X-1<Z+\mathcal M/X<\frac18.
\end{equation}
The opposite constant path was not discarded: its smallness follows
only after using equality and the integral minimum normalization.

\subsection{Balance becomes a gap in an integer exponent}

Define
\begin{equation}\label{eq:spectral-phase-normalization}
 q=m-\tau,\qquad
 n=\begin{cases}h-1,&a=c_0,\\h+1,&a=d_0.
 \end{cases}
\end{equation}
This internal offset may differ from $m-t$, because the dominant-state
choice can have $\gamma=-1$. The original enumeration convention of
\eqref{eq:gap-phase-window} is unchanged.
As in \eqref{eq:positive-coefficient-normalization}, exact rearrangement
gives
\begin{equation}\label{eq:spectral-coefficient-normalization}
 \begin{gathered}
 c_0=\frac{\lambda^{q/n}}3\alpha,\qquad
 d_0=\frac{\lambda^{-q/n}}3\beta,\qquad \beta=\kappa/\alpha,\\
 \alpha=\begin{cases}
 (1+\delta)^{-1/n},&a=c_0,\\
 [\kappa/(1+\delta)]^{1/n},&a=d_0.
 \end{cases}
 \end{gathered}
\end{equation}
Here $e_0=4/D<1/8$, but there is no assumption that $\delta<e_0$.
Positive-root monotonicity gives
\begin{equation}\label{eq:spectral-coefficient-bounds}
 \begin{gathered}
 8/9<\alpha<9/8,\qquad 1<\beta<81/64,\\
 |\alpha-1|,|\beta-1|<2(e_0+\delta),\qquad
 \max(\alpha/\beta,\beta/\alpha)<3/2.
 \end{gathered}
\end{equation}
For completeness, $\alpha$ lies between $1/(1+\delta)$ and $\kappa$,
and $1<\beta\leqslant\kappa(1+\delta)$. These bounds imply the
displayed error estimates; the ratio bound follows from
$(81/64)/(8/9)=729/512<3/2$.
If $q=0$, then $u_0=(\alpha+\beta)/3<51/64<1$, impossible.

Set
\begin{equation}\label{eq:spectral-gcd-scale}
 s=\lambda^{1/n}>S,\qquad j_0=|q|>0,
 \qquad g=\gcd(n,j_0).
\end{equation}
The minimum ratio and \eqref{eq:spectral-coefficient-bounds} give
$s^{2j_0-n}<3/2$ whenever $2j_0-n$ is positive. Since $s>4$,
integrality implies $1\leqslant j_0\leqslant n/2$.
As in the preceding section, $n=1$ is thereby excluded immediately.
The remaining cases have $n\geqslant2$, and
$2h/n\leqslant3$: this is immediate for $h\geqslant3$; at $h=2$
the only remaining choice is $n=3$.

For each proper lower block define the integers
\begin{equation}\label{eq:spectral-integer-activities}
 F_i=n+qj_i+\gamma ne_i,
 \qquad E=hq+n\tau=hm+(n-h)\tau.
\end{equation}
Since $m\geqslant1$, $n-h=\pm1$ and $\tau=\pm t$,
$E\geqslant h-t>0$. Moreover
\begin{equation}\label{eq:spectral-activity-positive}
 F_i=n\{1+\gamma(e_i-tj_i/h)\}+Ej_i/h>0.
\end{equation}
Both $E$ and every $F_i$ are multiples of $g$, so $E\geqslant g$
and $\sum_iF_i\geqslant rg$. This is where the actual equality
improves the coarse denominator $h$ in
\eqref{eq:spectral-coarse-block-weight} to the integer divisor $g$.

Put $R=\beta/\alpha>1$ and $a_g=\sigma\lambda^2s^{-2g}$.
Rewriting the exact weight gives
\begin{equation}\label{eq:spectral-normalized-block-weight}
 w=(\sigma\lambda^2)^rR^j s^{-2\sum_iF_i}
 \leqslant a_g^rR^j,
 \qquad Z=R^h s^{-2E}.
\end{equation}
The coefficient normalization also gives
$$
 R=\kappa^{\zeta}(1+\delta)^{2/n},\qquad
 \zeta=\begin{cases}1,&a=c_0,\\1-2/n,&a=d_0.
 \end{cases}
$$
In the case $a=d_0$ one has $n=h+1\geqslant3$ by
\eqref{eq:spectral-phase-normalization}, so $1/3\leqslant\zeta\leqslant1$ in both
cases. Since
$he_0\leqslant4h/(5p^2)<1/2$, the binomial/geometric estimate yields
$\kappa^h<1/(1-he_0)<2$. Together with $2h/n\leqslant3$ this gives
\begin{equation}\label{eq:spectral-ratio-power}
 R^h<2(9/8)^3<3.
\end{equation}
Equations \eqref{eq:spectral-partition-bound},
\eqref{eq:spectral-normalized-block-weight} and
\eqref{eq:spectral-ratio-power}, with
$h^2a_g<4h^2/S^2<1/16$, imply
\begin{equation}\label{eq:spectral-refined-delta}
 0<\delta<3s^{-2g}+3\Psi_h(a_g)
 <\left(3+\frac{64}{5}h^2\right)s^{-2g}.
\end{equation}
Also $e_0=4s^{-2n}/(1-s^{-2n})^2<5s^{-2n}$. Thus
\begin{equation}\label{eq:spectral-total-coefficient-error}
 e_0+\delta<\left(8+\frac{64}{5}h^2\right)s^{-2g}
 \leqslant15h^2s^{-2g},
\end{equation}
where the last inequality holds for every $h\geqslant2$.

Take the actual ray neighbor of $u_0$ at index
$-\operatorname{sign}(q)$ in the transformed ray. These two entries
and $p$ form a positive integral Markov triple with baseline indices
$(j_0,n-j_0,n)$. Using the scalar baselines
\eqref{eq:positive-baselines}, the exact coefficients give
\begin{equation}\label{eq:spectral-initial-label-errors}
 L_n=p=B_n,\qquad
 |L_i-B_i|<20h^2s^{i-2g},\qquad i=j_0,n-j_0.
\end{equation}
Unlike \eqref{eq:positive-initial-errors}, this error is measured at
the fixed divisor scale $g$, rather than at the initial maximum index
$n$. Direct propagation of relative errors would therefore lose a
factor at each mutation. Spectral coordinates avoid that loss.

\subsection{The exact spectral error of a Vieta mutation}

For a real label $L>2/3$, define
\begin{equation}\label{eq:spectral-angle}
 \theta(L)=\arcosh(3L/2).
\end{equation}
Then $\theta(B_i)=i\log s$ for $i>0$.
Let $L_a,L_b,L_c$ be a positive integral Markov triple, with $L_c$
strictly largest, and put $A=\theta(L_a)$, $B'=\theta(L_b)$.
Solving the quadratic in the third full trace gives
\begin{equation}\label{eq:spectral-large-root}
 \frac32L_c=\cosh A\cosh B'
       +\sqrt{\sinh^2A\sinh^2B'-1}.
\end{equation}
The Vieta partner $L_d=3L_aL_b-L_c$ is therefore determined by
\begin{equation}\label{eq:spectral-small-root}
 \begin{aligned}
 \frac32L_d&=\cosh(A-B')+\epsilon(A,B'),\\
 \epsilon(A,B')&=
 \frac1{\sinh A\sinh B'+\sqrt{\sinh^2A\sinh^2B'-1}}.
 \end{aligned}
\end{equation}
The square root is real and positive: $L_a,L_b\geqslant1$ give
$\sinh^2A,\sinh^2B'\geqslant5/4$. The sign in
\eqref{eq:spectral-large-root} is forced by the largest-entry
hypothesis. Indeed, if $x\geqslant y\geqslant3$ are the other
full traces, then the trace quadratic evaluated at $z=x$ is
$$
 x^2(2-y)+y^2\leqslant y^2(3-y)\leqslant0.
$$
A third trace strictly greater than $x$ must be its larger root.
The expression for $\epsilon$ follows by rationalizing the difference
$\sinh A\sinh B'-\sqrt{\sinh^2A\sinh^2B'-1}$.

At a nonterminal baseline stage let $a>b\geqslant g$,
$c=a+b$, and $d=a-b$. Suppose all three angle errors from their
baselines are less than $1/8$. Since each distinct-index gap is at
least $g\log s$, the actual entries have the same strict ordering as
their baseline indices. In particular \eqref{eq:spectral-large-root}
uses the correct root, and $A>B'$.
For $i\geqslant g$, $\abs v\leqslant1/4$ and $s\geqslant4$,
\begin{equation}\label{eq:spectral-sinh-bound}
 \sinh(i\log s-v)
 \geqslant\frac38s^i(1-2s^{-2i})
 \geqslant\frac{21}{64}s^i>\frac14s^i.
\end{equation}
Here $\exp(-1/4)>3/4$ and $\exp(1/2)<2$, as follows from their
power series. Apply this estimate to $A$, $B'$ and $A-B'$. The
inverse-cosh mean-value estimate in \eqref{eq:spectral-small-root}
then gives an exact additive error relation
\begin{equation}\label{eq:spectral-additive-forcing}
 \theta(L_d)=A-B'+\eta,
 \qquad 0<\eta<64s^{-2a}\leqslant64s^{-4g}.
\end{equation}
Indeed $\eta\leqslant\epsilon(A,B')/\sinh(A-B')$, and the three
sinh factors are bounded below by
$s^a/4,s^b/4,s^{a-b}/4$. Equation~\eqref{eq:spectral-additive-forcing}
is the reason to use spectral rather than relative label errors.

\subsection{Euclidean rows give a polynomial error bound}

\begin{lemma}[Coefficients along subtractive Euclid]
\label{lem:euclidean-row-bound}
Start with positive integer indices $a,b$ of gcd $g$, and put
$N=(a+b)/g$. Under repeated subtraction of the smaller index from
the larger, every signed coefficient row expressing a current index
in terms of the starting pair has $\ell^1$ norm at most $N$.
The same bound applies to propagation from any intermediate stage.
If the two initial coordinate errors are bounded by $e_*$ and every
additive update error is bounded by $f_*$, then all later coordinate
errors are bounded by
\begin{equation}\label{eq:spectral-row-error-bound}
 Ne_*+N(N-1)f_*<Ne_*+N^2f_*.
\end{equation}
\end{lemma}

\begin{proof}
Track the two coefficient rows $r,s$, initially $(1,0),(0,1)$.
They remain unimodular. In each coordinate $i$, the weak opposite-sign
condition $r_i s_i\leqslant0$ is preserved: replacing $r$ by $r-s$
changes the product to $r_i s_i-s_i^2\leqslant0$.
Consequently $|r-s|=|r|+|s|$ componentwise. Each new absolute row
dominates both rows that produced it.

Include the final subtraction of equal indices, producing zero and
$g$. If $A_*=a/g$ and $B_*=b/g$, the zero row is a primitive
integer relation on the coprime pair $(A_*,B_*)$, hence is
$\pm(B_*,-A_*)$. It dominates every preceding absolute row and has
$\ell^1$ norm $N$. This proves the row bound.
At any intermediate stage the gcd is still $g$, while the normalized
sum is at most $N$, so the same proof applies to every remaining tail.
There are at most $N-1$ subtractions, since each decreases the normalized
sum by at least one. Propagating each additive error along its remaining
tail and summing proves \eqref{eq:spectral-row-error-bound}.
Coordinate exchanges have norm one and require no additional estimate.
\end{proof}

We now apply the lemma to the angle errors, with initial nonlargest
indices $j_0,n-j_0$. The actual index-$n$ entry has zero initial
error. From \eqref{eq:spectral-initial-label-errors}, every point on
the segment between $L_i$ and $B_i$ is at least $s^i/4$, because
$20h^2s^{-2g}<1/12$. On this segment
$$
 \theta'(L)=\frac1{\sqrt{L^2-4/9}}<8s^{-i}.
$$
The initial angle errors are therefore bounded by
$e_*=160h^2s^{-2g}$. We have $N=n/g\leqslant h+1\leqslant2h$.
If the errors stay below $1/8$, equation~\eqref{eq:spectral-additive-forcing}
permits $f_*=64s^{-4g}$. The bound from
Lemma~\ref{lem:euclidean-row-bound} is then
\begin{equation}\label{eq:spectral-bootstrap-margin}
 Ne_*+N^2f_*
 <\frac{320h^3}{S^2}+\frac{256h^2}{S^4}
 \leqslant\frac2{125}+\frac{256}{10^8}<\frac1{16}.
\end{equation}
The constant $2/125$ includes the boundary $h=2$.

This proves the assumed error bound by a first-violation argument.
Initially all errors are below $1/8$. Before a first violating move,
the previously proved ordering selects the larger root at every stage,
and every forcing is bounded by \eqref{eq:spectral-additive-forcing}.
The row estimate then bounds the new error by
\eqref{eq:spectral-bootstrap-margin}, a contradiction. Only the two
nonlargest baseline coordinates are updated as a subtractive pair.
The third entry is initially exact, and thereafter is an unchanged
entry from the previous pair, so it has the same error bound. There is
no extra independent error from a discarded large coordinate.

\subsection{The terminal contradiction}

The subtractive Euclidean algorithm reaches $(g,g,2g)$.
The two index-$g$ angles then satisfy $|A-B'|<1/8$, while each is
at least $g\log s-1/16$. By \eqref{eq:spectral-sinh-bound},
$\epsilon(A,B')<16s^{-2g}<1/16$.
The actual index-$2g$ entry is still strictly largest, so its positive
integral Vieta partner satisfies \eqref{eq:spectral-small-root}.
Using $\cosh(1/4)<\sum_{j\geqslant0}(1/4)^{2j}=16/15$, we obtain
$$
 0<L_d<\frac23\left(\frac{16}{15}+\frac1{16}\right)
       =\frac{271}{360}<1.
$$
This contradicts positive integrality. It includes the equal-index
starting case $j_0=n/2$, with no nonterminal steps. Together with the
earlier direct exclusion of $n=1$, it proves
Theorem~\ref{thm:minimum-bound}.

Two different estimates were needed. The first used cyclic balance
before any approximation assumption to select a small discrepancy.
The second used the exact equality to make all block exponents
positive multiples of $g$. Spectral subtraction then transported the
resulting error with polynomial coefficients. The theorem supplies a
finite parameter range at every prescribed occurrence count; it does
not assert that a constant-state path dominates uniformly when $h$
grows on a fixed ray. The next section turns the two parameter bounds
into a complete finite comparison and isometry decision.
